\documentclass{article}

\usepackage{fontspec}
\usepackage{polyglossia}
\setmainlanguage{russian}
\setotherlanguage{english}
\defaultfontfeatures{Ligatures=TeX}

\setmainfont{DejaVuSerif.ttf}[
  BoldFont       = DejaVuSerif-Bold.ttf,
  ItalicFont     = DejaVuSerif-Italic.ttf,
  BoldItalicFont = DejaVuSerif-BoldItalic.ttf,
]
\setmonofont{DejaVuSansMono.ttf}[
  BoldFont       = DejaVuSansMono-Bold.ttf,
  ItalicFont     = DejaVuSansMono-Oblique.ttf,
  BoldItalicFont = DejaVuSansMono-BoldOblique.ttf,
]

\usepackage{amsmath,amssymb,amsfonts}
\usepackage{graphicx}
\usepackage{booktabs}
\usepackage{longtable}
\usepackage{array}
\usepackage{hyperref}
\usepackage{url}
\usepackage{float}
\usepackage{tabularx}


\begin{document}

\title{Dual Scale Inversion Cosmology (DSIC): фоновая космология одного параметра из относительности масштабов}

\author{Азамат Засеев\\
Independent Researcher, Russia\\
ORCID: 0009-0004-9345-1069\\
\url{https://doi.org/10.5281/zenodo.21198494}\\
\url{https://github.com/az-zase/DSIC}}

\date{}
\maketitle

\begin{abstract}
Представлена геометрическая модель космологического фона --- Dual Scale Inversion Cosmology (DSIC), --- в которой пространство и объектность рассматриваются как два полюса одной сохраняющейся величины, а абсолютный масштаб отсутствует: физическим содержанием обладает лишь отношение полюсов. Состояние системы движется по окружности $S^2 + O^2 = R^2$ под действием единственного равномерного потока; внутренний наблюдатель, сам построенный из объектности, имеет доступ только к проекции. Из линейки и часов такого наблюдателя однозначно выводится наблюдаемый масштабный фактор $b(\mu) = \sqrt{1-\mu^2}/\mu^2$, а из него --- красное смещение, кинематика и расстояния. Ускорение поздней Вселенной оказывается кинематическим следствием относительного разбега масштабов двух полюсов; тёмная энергия и аналоги $\Omega_m$, $\Omega_\Lambda$ в модель не вводятся.

Единственный параметр $\mu_0 = 0.7600 \pm 0.0091$ (фаза наблюдателя «сегодня»; $1\sigma$ из гессиана $\chi^2$) согласованно описывает шесть независимых проб фоновой кинематики при $z < 2.3$. Форма зависимости $b(\mu)$ фиксирована конструкцией и не содержит функциональной свободы, поэтому измеренное поздними данными значение $\mu_0$ однозначно определяет поведение модели на высоком $z$: расхождения с $\Lambda$CDM при $z > 3$ являются вынужденными предсказаниями, а не результатом подстройки. На полном каталоге Pantheon+SH0ES [3, 4] с ковариацией STAT+SYS модель статистически неотличима от плоской $\Lambda$CDM ($\Delta\chi^2 \approx +0.05$ при равном числе свободных параметров); совместный фит SN+BAO формально слабо предпочитает DSIC ($\Delta\chi^2 \approx -5.0$, что частично отражает известное натяжение DESI--$\Lambda$CDM). $Om(z)$-диагностика [11] на текущих данных моделей не различает (null-результат: знак $\Delta\chi^2$ зависит от нормировки $H_0$). Наиболее жёсткие проверки ядра: радиальный масштаб BAO в области Lyman-$\alpha$ (расхождение с $\Lambda$CDM до ${\sim}5\%$ при $z \approx 3.5$, досягаемо для DESI DR3) и бюджет космического времени при $z = 5\text{--}12$ (при $z = 7.5$ доступно 0.42 Гыр против 0.71 у $\Lambda$CDM), количественно разобранные в \S11.6.

Ранняя Вселенная вынесена во \textbf{второй этаж} модели (Часть II): форма поправки к пути света выводится из дополнительного постулата о связанности объектов, а единственный порог $\mu_d \approx 0.9806$ ($z \approx 4.5$) калибруется по одному наблюдению --- расстоянию до поверхности последнего рассеяния; статус этой надстройки --- калибровка при фиксированной постулатом форме, а не предсказание --- оговаривается явно. Рост структур не выводится из ядра; фон модели проверен на совместимость со стандартным (GR-предельным) ростом линейных возмущений: на компиляции $f\sigma_8$ («Gold-2017» [13, 14]) восстанавливается $S_8 \approx 0.76$, статистически неотличимо от $\Lambda$CDM (\S12); вывод роста из связанности остаётся направлением развития. Все эмпирические утверждения воспроизводятся скриптами из открытого репозитория.
\end{abstract}
\section{Введение}

Переход космологического расширения от замедления к ускорению при $z_t \approx 0.6\text{--}0.8$ [1, 2] --- твёрдо установленный наблюдательный факт. Стандартная модель описывает его введением компоненты с отрицательным давлением: два параметра плотности $\Omega_m$ и $\Omega_\Lambda$ задают форму кривой расширения, но сама постоянная $\Lambda$ остаётся постулатом без вывода. Естественен вопрос: может ли та же кривая расширения возникнуть без энергетического ингредиента --- из одной лишь геометрии того, что мы называем «масштабом»?

DSIC даёт на этот вопрос утвердительный ответ в форме явной конструкции. Модель исходит из того, что абсолютного масштаба не существует ни для объектов, ни для пространства: наблюдаемо только их отношение. Пространство и объектность трактуются как два полюса одной сохраняющейся нормы, перетекающей по окружности под действием единственного равномерного потока. Наблюдатель, находящийся внутри системы и сам состоящий из объектности, принципиально не имеет доступа ни к знаку величин, ни к их абсолютным значениям --- только к модулям и отношениям. Требование постоянства скорости света в его единицах однозначно фиксирует его часы, и из линейки и часов выводится наблюдаемый масштабный фактор
\begin{equation}
b(\mu) = \frac{\sqrt{1 - \mu^2}}{\mu^2}\,,
\end{equation}
где $\mu \in [0, 1]$ --- наблюдаемая доля объектности. Эта единственная функция порождает красное смещение, темп расширения, параметр замедления и все расстояния. Смена знака параметра замедления --- переход от торможения к ускорению --- встроена в кривизну $b(\mu)$ и не требует ни $\Lambda$, ни какой-либо новой компоненты: ускорение возникает как кинематическое следствие нарастающего разбега двух масштабных полюсов.

Заявленная область применимости оговаривается сразу. Ядром показана \textbf{фоновая кинематика поздней Вселенной} ($z < 2.3$): один параметр $\mu_0$ описывает шесть независимых проб статистически не хуже плоской $\Lambda$CDM при равном числе свободных параметров. Ключевое свойство модели --- не экономность, а жёсткость: форма $b(\mu)$ фиксирована конструкцией и не содержит функциональной свободы, поэтому значение $\mu_0$, измеренное поздними данными с точностью ${\sim}1\%$, однозначно определяет поведение модели на высоком $z$ --- включая $E(z)$, радиальный масштаб BAO в области Lyman-$\alpha$ и бюджет космического времени при $z = 5\text{--}12$ --- без возможности подстройки. Эти предсказания расходятся с $\Lambda$CDM на уровне, досягаемом для DESI DR3, и составляют прямой тест ядра (\S11.6). Ранняя Вселенная дотягивается \textbf{вторым этажом} модели (Часть II): одна новая сущность --- связанность объектов --- фиксирует постулатом форму поправки к пути света, а единственный порог $\mu_d$ калибруется по одному наблюдению --- расстоянию до реликта; статус этой надстройки (калибровка, а не предсказание) и её ограничения разобраны в \S11.6. Рост структур в ядро не входит; фон модели проверен лишь на совместимость со стандартным ростом линейных возмущений (\S12).

Изложение подчинено строгой слоевой дисциплине, и по виду символа всегда видно, к какому слою он принадлежит:
\begin{itemize}
  \item \textbf{слой 0 --- поток}: единственный монотонный параметр $\varphi$ (время);
  \item \textbf{слой 1 --- ядро}: величины со знаком $S, O, \sigma$; здесь нет ни сил, ни циклов, ни разворотов, и операция модуля не применяется;
  \item \textbf{слой 2 --- проекция}: величины без знака $P, Q, \mu$, доступные внутреннему наблюдателю; рождаются единственной операцией $|\cdot|$; отскоки и циклы существуют только здесь;
  \item \textbf{слой 3 --- наблюдаемая физика}: $b, z, H_\tau, q, \chi, d_L$ --- интерпретация проекции в физических единицах через шкалы наблюдателя.
\end{itemize}

Поскольку модель строится на собственном языке, приведём словарь соответствий между её терминами и привычными физическими понятиями. Термины намеренно сохранены (они несут онтологию модели), но каждому дан стандартный аналог:

\begin{table}[htbp]
\centering
\small
\begin{tabular}{@{}p{2.4cm}p{3.2cm}p{3.2cm}@{}}
\toprule
\textbf{термин DSIC} & \textbf{что это} & \textbf{ближайший привычный образ} \\
\midrule
объектность ($O$) & полюс «вещественности», масштаб объектов & плотность материи / вещество \\
\addlinespace
пространство ($S$) & полюс «протяжённости», масштаб пространства & объём / пространственный масштаб \\
\addlinespace
норма $R$ & сохраняющаяся сумма полюсов $S^2 + O^2 = R^2$ & закон сохранения (связь двух полюсов) \\
\addlinespace
поток $\varphi$ & единственная монотонная степень свободы & время как параметр эволюции \\
\addlinespace
каша & предел $\mu \to 1$: максимум объектности, минимум пространства & регулярный аналог, близкий к начальной сингулярности \\
\addlinespace
отскок объектов & предел $\mu \to 0$: максимум пространства & асимптотика де Ситтера \\
\addlinespace
выворот & смена знака фазы в ядре (проход полюсом нуля) & зеркальная фаза, ненаблюдаемая изнутри \\
\addlinespace
отскок & наблюдаемый разворот у предела & артефакт проекции (взятия модуля), а не событие \\
\addlinespace
проекция ($\mu$, $P$, $Q$) & модульные, наблюдаемые величины & то, что доступно измерению \\
\addlinespace
связанность (Часть II) & мера сцепления объектов & аналог гравитационной связанности \\
\addlinespace
отлипание ($\mu_d$) & порог, за которым монолитная материя позволяет проступить пространству достаточно, чтобы свет свободно распространялся & эпоха, после которой геометрия ядра работает без поправки \\
\bottomrule
\end{tabular}
\end{table}

Ключевое, что стоит держать в уме: «отскок» и «каша» --- не физические события, а то, как внутренний наблюдатель видит гладкое монотонное движение ядра через операцию модуля (\S4.3). Под наблюдаемым колебанием между двумя пределами лежит одно движение вперёд, в котором ничто не возвращается.

Граница между слоями абсолютна: знак из ядра в проекцию не проходит (модуль его срезает), и ни одна величина проекции не входит в уравнения ядра.

План статьи. Часть I --- ядро: \S2 формулирует постулаты; \S3 строит ядро (окружность, инверсии, знаковые фазы); \S4 --- проекцию для внутреннего наблюдателя; \S5 выводит наблюдаемую физику; \S6 сверяет модель с данными; \S7 перечисляет фальсифицируемые предсказания; \S8 очерчивает границы; \S9 позиционирует модель относительно $\Lambda$CDM. Часть II --- второй этаж: \S10 вводит связанность объектов и постулат отлипания; \S11 строит раннюю Вселенную (порог $\mu_d$); Приложение B --- модуль распределения пространства (гравитация). \S13 --- общее заключение. Приложение A содержит сводку обозначений.

\bigskip
\hrule
\bigskip

\part{Ядро и наблюдаемая космология}
\section{Постулаты}

\noindent\textbf{P1 (относительность масштаба).} Абсолютного масштаба не существует. Физическим содержанием обладает только отношение масштабов; любая конструкция, открывающая доступ к абсолютной шкале, запрещена.

\medskip
\noindent\textbf{P2 (два полюса, одна норма).} Объектность и пространственность --- не две независимые сущности, а два полюса одного отношения, связанные общей степенью свободы. Объект --- локальная выпуклость: имеет внутреннюю область и внешнюю границу, занимает место относительно других объектов (шар, камень, планета, атом --- варианты одной структуры). Пространство --- антипод объекта, всеобщая вогнутость: оно не находится нигде и не содержится ни в чём; сама категория «где» возникает только внутри него. Масштабы полюсов $S$ (пространство) и $O$ (объектность) подчинены закону сохранения нормы:
\begin{equation}
S^2 + O^2 = R^2 = \text{const}\,.
\end{equation}
Рост одного полюса возможен лишь за счёт другого; физику несёт фиксированность нормы при перетекании, а не абсолютные значения полюсов.

\medskip
\noindent\textbf{P3 (единственный поток).} У системы одна фундаментальная степень свободы --- непрерывный монотонный поток $\varphi$, никогда не разворачивающийся. Его динамика свободна: действие не содержит ни потенциала, ни сил, ни космологической постоянной,
\begin{equation}
\mathcal{S}[\varphi] = \int dt\, \tfrac{1}{2}\dot{\varphi}^2\,, \qquad \delta\mathcal{S} = 0 \;\Rightarrow\; \ddot{\varphi} = 0 \;\Rightarrow\; \varphi(t) = \omega\,t\,, \quad \omega > 0\,.
\end{equation}
Особых точек у $\varphi$ нет; всё остальное в модели --- функции от $\varphi$. Цикличность, если она наблюдается, рождается не из потока, а из способа его наблюдать.

\medskip
\noindent\textbf{P4 (внутренний наблюдатель).} Любой наблюдатель находится внутри системы и сам является объектом: его приборы преобразуются вместе с ней. Ему недоступны знак величин и абсолютный масштаб; доступны только модули --- на практике лишь сближение и отдаление объектов.

\medskip
\noindent\textbf{P5 (постоянство скорости света).} В единицах внутреннего наблюдателя $c = \text{const}$. Как показано в \S5.1, вместе с P1 это однозначно фиксирует ход его часов.
(Это фиксированный конвенцией курс, связывающий линейку и часы наблюдателя --- то есть ещё одно отношение, а не абсолютная величина. Численное значение $c$ нигде не входит в вывод геометрии и появляется лишь при переводе безразмерных предсказаний в единицы наблюдений.)

\section{Ядро: окружность, инверсии, знаковые фазы}

\subsection{Несущая и окружность}

Оба полюса строятся из одной гладкой несущей, сдвинутой на четверть оборота (квадратура):
\begin{equation}
S(\varphi) = R\sin(\varphi/2)\,, \qquad O(\varphi) = R\cos(\varphi/2)\,, \qquad S^2 + O^2 = R^2\,.
\end{equation}
Несущая имеет период $4\pi$ по фазе потока --- полный знаковый цикл. Полное описание обладает определённой чётностью и постоянным вронскианом:
\begin{align}
S(-\varphi) &= -S(\varphi)\,, \qquad O(-\varphi) = O(\varphi)\,, \\
W &= S\cdot\frac{dO}{d\varphi} - O\cdot\frac{dS}{d\varphi} = -\frac{R^2}{2} = \text{const}\,.
\end{align}
Постоянство $W$ --- формальная запись квадратуры: когда один полюс в экстремуме, другой проходит ноль; всякая разрывная производная модуля (возникающая позже, в проекции) умножается на обращающийся в ноль парный множитель, так что произведение в полном описании со знаком остаётся гладким.

Если отложить $S$ и $O$ по двум осям, инвариант рисует окружность радиуса $R$: всякое допустимое состояние системы --- точка на ней. Эволюция --- скольжение этой точки с постоянной угловой скоростью (равномерный поток $\varphi$); уйти внутрь или наружу окружности нельзя (норма фиксирована), а проходя точку на оси, состояние не останавливается и не отскакивает --- продолжает движение в ту же сторону. Скорость перетекания между полюсами при этом непостоянна:
\begin{equation}
\frac{d\mu}{d\varphi} \propto \sin(\varphi/2) \quad \text{--- медленно у полюсов, быстро в середине дуги.}
\end{equation}
Эта нелинейность --- чисто геометрическое свойство проекции равномерного кругового движения на оси, а не действие силы: сам поток остаётся строго равномерным (P3).

\subsection{Инверсии}

По окружности строго чередуются два события --- гладкие проходы нулей полюсов:
\begin{align}
\text{инверсия пространства:} &\quad S(\varphi) = 0 \;\Rightarrow\; \varphi = 2k\pi \qquad (k \in \mathbb{Z})\,, \\
\text{инверсия объектов:} &\quad O(\varphi) = 0 \;\Rightarrow\; \varphi = \pi + 2k\pi\,.
\end{align}
Нуль $S$ совпадает с максимумом $|O| = R$: это состояние \textbf{«каши»} --- предельной плотности объектности, в котором пространства почти не остаётся, а объекты сжаты до предела. Важно различать: каша --- вся зона вблизи нуля пространства, \textbf{с обеих сторон} от него; половина до нуля --- финал сжатия, половина после --- начало разлёта. \textbf{Большим Взрывом} уместно называть только вторую, исходящую половину, где модуль пространства снова растёт и объектность высвобождается. Один и тот же ноль пространства разделяет финал предыдущего сжатия и начало нового разлёта.

Нуль $O$ совпадает с максимумом $|S| = R$: это минимум объектности --- \textbf{отскок объектов} в терминах наблюдателя (см.\ \S4).

Из-за нелинейности перетекания (\S3.1) интервалы между инверсиями по $\varphi$ неравномерны, но порядок чередования строг: каша $\to$ отскок объектов $\to$ следующая каша.

\subsection{Знаковые фазы и наследование}

Пройдя свой ноль, полюс переходит в \textbf{зеркальную (знаковую) фазу}: продолжает то же движение с противоположным знаком --- выворот координат наизнанку, гладкий на уровне несущей. Обозначим знаковые фазы $\sigma_S = \operatorname{sign} S(\varphi)$, $\sigma_O = \operatorname{sign} O(\varphi)$.

Ключевая структурная асимметрия модели: объектность \textbf{дочерняя} по отношению к пространству, и два нуля устроены по-разному.
\begin{itemize}
  \item При $O = 0$ (отскок объектов) свой ноль проходит \textbf{только объектность} --- она и выворачивается, меняя знак фазы. Пространство в этот момент на максимуме и совершает лишь наблюдаемый отскок: его масштаб перестаёт прибывать и начинает убывать, но знак фазы пространства не меняется.
  \item При $S = 0$ (каша) свой ноль проходит \textbf{само пространство} --- выворачивается общий контейнер, и содержимое выворачивается вместе с ним: объектность, находясь в максимуме и не проходя собственного нуля, \textbf{наследует} инверсию знака принудительно, синхронно с родителем --- включая измерительные приборы наблюдателя.
\end{itemize}

События не зеркальны: при нуле объектности знак меняет один полюс, при нуле пространства --- оба (один сам, другой по наследству). Формально наследование выражается величиной
\begin{equation}
O_{\text{full}}(\varphi) = O(\varphi)\cdot\sigma_S(\varphi)\,,
\end{equation}
которая непрерывна в точках инверсии объектов ($O$ проходит ноль плавно, $\sigma_S$ постоянен) и разрывна в точках инверсии пространства (мгновенный переброс знака). Разрыв ненаблюдаем изнутри: $|O_{\text{full}}| = |O|$ непрерывна везде, а синхронная смена знака у системы и приборов сокращается. Именно непрерывность/разрывность $O_{\text{full}}$ определит в \S5.3 наблюдаемую асимметрию двух пределов масштаба.

Терминологическое соглашение, которого статья придерживается всюду: \textbf{выворот} --- смена знака фазы на уровне ядра (случается с полюсом, проходящим свой ноль); \textbf{отскок} --- то, что видит внутренний наблюдатель: касание предела и разворот. Это не синонимы, а описания одного механизма с двух позиций.

\begin{table}[htbp]
\centering
\small
\begin{tabular}{@{}p{2.6cm}p{3.6cm}p{3.6cm}p{2.2cm}@{}}
\toprule
\textbf{ноль} & \textbf{что выворачивается (меняет знак фазы)} & \textbf{второй полюс} & \textbf{излом «V» у наблюдателя} \\
\midrule
$O = 0$ (отскок объектов) & сама объектность & пространство на пике: просто отскок, знак неизменен & на $Q = |O|$ \\
\addlinespace
$S = 0$ (каша) & пространство $+$ объектность по наследству & объектность на пике, вывернута принудительно & на $P = |S|$ \\
\bottomrule
\end{tabular}
\end{table}

\subsection{Падение в себя}

Всеобщая вогнутость пространства (P2) задаёт и характер «схлопывания» объектности: падение в этой картине --- не движение к точке в пространстве, а схождение масштаба объекта к собственному центру, идущее внутрь отовсюду сразу, без выделенного направления. Объекты не летят в общую яму «где-то там» --- они проваливаются в себя. Гравитационное развитие этой мысли (общий центр падения у связанных групп) обсуждается в \S8 как направление, не выведенное из ядра.
\section{Проекция: мир внутреннего наблюдателя}

\subsection{Два наблюдателя}

Одно и то же событие выглядит по-разному с двух позиций. \textbf{Внешний наблюдатель} --- воображаемая точка зрения вне пространства (в действительности её нет: пространство фундаментально и не имеет «снаружи», но она удобна для полного описания). Он видит поток целиком, со знаком: история непрерывна и монотонна, величины проходят нули и продолжают меняться в ту же сторону; никаких разворотов. \textbf{Внутренний наблюдатель} --- любой, кто внутри. По P4 ему доступны только модули: там, где внешний видит гладкий проход нуля, внутренний видит касание нуля и разворот --- отскок. И тихий выворот наизнанку, и гладкий проход нуля читаются изнутри одинаково --- как смена режима сближения/отдаления объектов, наблюдаемая существом, которое само сделано из того, что меняется.

\subsection{Операция модуля}

Переход от ядра к наблюдаемому осуществляет единственная операция --- взятие модуля:
\begin{equation}
\pi : (S, O) \mapsto (|S|, |O|)\,.
\end{equation}
Она делает три вещи: срезает знак (зеркальная ветвь недоступна); сворачивает гладкий проход нуля в излом --- отскок; и отождествляет соседние знаковые ветви, отчего наблюдаемый период оказывается вдвое короче знакового. Наблюдаемые величины:
\begin{align}
P(\varphi) &= |S| = R\,|\sin(\varphi/2)| \ge 0\,, \qquad Q(\varphi) = |O| = R\,|\cos(\varphi/2)| \ge 0\,, \\
&\quad P^2 + Q^2 = R^2\,, \notag \\[4pt]
\mu(\varphi) &= Q/R = |\cos(\varphi/2)| \in [0, 1]\,, \qquad P = R\sqrt{1 - \mu^2}\,, \quad Q = R\mu\,.
\end{align}
Наблюдаемая фаза масштаба $\mu$ (доля объектности) имеет период $2\pi$ против знакового $4\pi$: два различных в полном описании события сливаются для наблюдателя в одно.
\begin{align*}
\text{полное описание (со знаком):} &\quad \ldots\; +O \to 0 \to -O \;\ldots \quad (\text{монотонный проход}) \\
\text{наблюдаемое (модуль):} &\quad \ldots\; |O| \to 0 \to |O| \;\ldots \quad (\text{отскок})
\end{align*}
Пара $(P, Q)$ неотрицательна и всегда лежит в первой четверти. За один наблюдаемый цикл точка проходит эту четверть туда и обратно: от каши ($P = 0$, $Q = R$) до отскока объектов ($P = R$, $Q = 0$) и назад. Движение «туда-обратно» возникает не из наложения дуг, а из симметрии $Q = R\,|\cos(\varphi/2)|$ относительно точки отскока: дойдя до $Q = 0$, точка отражается и возвращается по той же дуге. Монотонное круговое скольжение ядра выглядит изнутри как колебание между двумя пределами; число знаковых тактов ядра при этом ненаблюдаемо --- это свойство границы слоёв, а не неопределённость модели.

\subsection{Отскок --- не физическое событие}

Несущая гладкая, собственных вершин у неё нет, поэтому \textbf{оба} наблюдаемых излома $\mu$ имеют одну природу --- их создаёт модуль при прохождении соответствующим полюсом нуля: при $O = 0$ в излом «V» складывается $Q = |O|$ (в этой точке $P$ гладко проходит максимум), при $S = 0$ --- $P = |S|$ (гладок максимум $Q$). Отскок не требует силы, торможения или причины: разворота на уровне механизма нет вовсе, есть смена способа видеть. «Проход через ноль» и «отскок» --- одно событие с двух позиций.

За один цикл внутренний наблюдатель видит четыре режима: (i) расхождение объектов доходит до предела ($\mu$ минимальна) --- начинается рост объектного масштаба; (ii) объекты сближаются, пространство убывает; (iii) каша --- объекты сталкиваются и приобретают сверхплотное пережатое пространством состояние; (iv) каша расходится (Большой Взрыв) --- объекты отдаляются. Все повороты выглядят однотипными сменами режима, но повторяется проекция, а не течение: под ней --- одно монотонное движение вперёд, в котором ничто не возвращается.
\section{Наблюдаемая физика}

Всё в этом разделе строится из $\mu$ и шкал внутреннего наблюдателя; это интерпретация проекции в физических единицах, а не отдельный механизм.

\subsection{Масштабный фактор}

Наблюдатель сделан из объектности, поэтому его единица длины стянута к объектному полюсу: линейка $L(\mu) \propto \mu$. Метрика --- отношение полюсов $P/Q = \sqrt{1 - \mu^2}/\mu$; у нуля объектности линейка стягивается в точку, и относительные измерения расходятся.

Ход часов фиксируется постулатом P5. Период собственных часов $T(\mu) \propto \mu^p$; тогда скорость света изнутри $c = L/T \propto \mu^{1-p}$. Требование $c = \text{const}$ даёт единственное значение $p = 1$, то есть собственное время наблюдателя
\begin{equation}
d\tau = dt / \mu\,.
\end{equation}
Любой другой показатель $p$ гнал бы $c$ по диапазону и открывал доступ к абсолютному масштабу, запрещённому P1. Наблюдаемый масштаб складывается из метрики и хода часов:
\begin{equation}
b(\mu) = \frac{P}{Q}\cdot\frac{1}{\mu} = \frac{\sqrt{1 - \mu^2}}{\mu^2}\,.
\end{equation}
Это центральная формула модели: одна функция одной переменной, без свободных параметров формы.

\subsection{Красное смещение}

Свет \textbf{распространяется по пространству} (полюс $P$), но \textbf{измеряется} наблюдателем, построенным из объектности ($L \propto \mu$, $d\tau = dt/\mu$, $c = \text{const}$). Распространение и измерение идут по разным полюсам, поэтому в наблюдаемую длину волны входят оба. На эпохе с фазой $\mu$ она складывается из трёх вкладов, каждый из которых входит ровно один раз:
\begin{align*}
&(1)\; \text{растяжение пути по пространству:} && \propto \sqrt{1 - \mu^2} && [\text{полюс } P] \\
&(2)\; \text{часы источника, } d\tau = dt/\mu: && \propto 1/\mu && [\text{время}] \\
&(3)\; \text{объектная линейка, } L \propto \mu: && \propto 1/\mu && [\text{длина}] \\[4pt]
&\phantom{(3)\;} \lambda_{\text{obs}}(\mu) \propto \sqrt{1 - \mu^2} / \mu^2 = b(\mu)\,.
\end{align*}
Наблюдаемая длина волны тождественно равна масштабному фактору, и красное смещение выпадает из счёта вкладов автоматически --- оно выведено, а не введено:
\begin{equation}
1 + z = \frac{\lambda_{\text{obs}}(\mu_0)}{\lambda_{\text{obs}}(\mu_e)} = \frac{b(\mu_0)}{b(\mu_e)} = \frac{\mu_e^2}{\mu_0^2}\cdot\frac{\sqrt{1 - \mu_0^2}}{\sqrt{1 - \mu_e^2}}\,,
\end{equation}
где $\mu_0$ --- фаза приёма («сегодня»), $\mu_e$ --- фаза излучения. При $\mu_e \to 1$ (каша) $z \to \infty$; при $\mu_e \to \mu_0$ $z \to 0$. Формула допускает разложение на два полюса:
\begin{equation}
1 + z = \underbrace{\frac{\sqrt{1 - \mu_0^2}}{\sqrt{1 - \mu_e^2}}}_{\text{полюс } P\ (\text{пространство})} \cdot \underbrace{\frac{\mu_e^2}{\mu_0^2}}_{\text{объектная шкала}}
\end{equation}
--- первый множитель отвечает растяжению пространства вдоль пути, второй --- объектной шкале наблюдателя (часы $+$ линейка). Их сменное доминирование вдоль цикла и есть механизм наблюдаемой кинематики (\S5.4). Отметим и физику предела: у каши полюс $P$ стягивается почти в ничто --- свету попросту почти не по чему распространяться; путь по пространству схлопывается, и наблюдаемая картина держится уже на объектной шкале. Свет не «замедляется» --- исчезает сцена, по которой он шёл.

Альтернатива «редшифт по одной линейке», $1 + z = \mu_0/\mu_e$, отвергается: она учитывает только вклад (3) и игнорирует путь (1) и часы (2), тогда как фотон проживает все три. Полная метрика наблюдателя есть $b$, поэтому редшифт идёт по $b$, а не по $L$.

Два свойства формулы. Во-первых, $b(\mu)$ строго убывает на $(0, 1)$, поэтому соответствие $z \leftrightarrow \mu_e$ на физической ветви приёма $\mu_e \in (\mu_0, 1)$ взаимно однозначно и инверсия $z \to \mu_e$ корректна. Во-вторых, в $1 + z$ входит только модульная величина $\mu$ --- знак из ядра не учитывается, слоевая граница не нарушена. Сам по себе $z$ в изоляции ненаблюдаем: в данные он входит только через $d_L(z)$.

\subsection{Пределы цикла и их асимметрия}

\begin{align*}
\mu \to 0 \;\; (\text{отскок объектов}): &\quad b \to \infty\,, \\
\mu \to 1 \;\; (\text{каша}): &\quad b \to 0\,.
\end{align*}
У отскока объектов пространство максимально, объекты исчезающе малы относительно него, и наблюдаемый масштаб расходится; это \textbf{полюс приближения, а не достигаемая сингулярность} --- наблюдатель скользит мимо $\mu = 0$, не попадая в неё. В каше отношения масштабов не остаётся вовсе. Асимметрия $\infty$ vs $0$ не противоречит симметрии ядра, а прямо следует из механики знака (\S3.3): через ноль объектности объектность проходит сама и гладко ($O_{\text{full}}$ непрерывна $\to$ расходимость), через ноль пространства проходит контейнер, выворачивая знак содержимого ($O_{\text{full}}$ разрывна $\to$ обнуление). Расходимость $b$ у нуля объектности --- прямое выражение относительности масштабов (ограниченные формы $b$ постулату P1 не удовлетворяют) и артефакт проекции, а не событие ядра.

\subsection{Кинематика}

Введём направление ветви $s = \pm 1$ ($s = -1$ --- расширение, $\mu$ убывает; $s = +1$ --- сжатие) и характеристическую скорость фазы $k = (2/\pi)\cdot\omega > 0$, так что вдоль ветви $d\mu/dt = s\cdot k = \text{const}$; множитель $2/\pi$ фиксирует масштаб времени так, чтобы наблюдаемый цикл по $\mu$ имел период $2\pi$. Все производные берутся по собственному времени $d\tau = dt/\mu$, то есть $d/d\tau = \mu\cdot d/dt$. Постоянство $k$ --- определение ветви, а не приближение: при непостоянном темпе в параметр замедления вошёл бы член $\propto dk/dt$, здесь тождественно равный нулю. Поэтому кинематика зависит только от фазы $\mu$, но не от того, как быстро система её проходит.

Темп расширения:
\begin{align}
H_\tau(\mu) &= \frac{db/d\tau}{b} = s\cdot k\cdot\frac{\mu^2 - 2}{1 - \mu^2}\,, \\
|H_\tau(\mu)| &= \frac{2 - \mu^2}{1 - \mu^2} \quad [\text{ед. } k]\,; \\
|H_\tau| &\to 2 \;\text{ при } \mu \to 0\,, \qquad |H_\tau| \to \infty \;\text{ при } \mu \to 1\,. \notag
\end{align}
Параметр замедления:
\begin{equation}
q(\mu) = -\frac{(d^2 b/d\tau^2)\cdot b}{(db/d\tau)^2} = \frac{-\mu^4 + 6\mu^2 - 4}{\mu^4 - 4\mu^2 + 4}\,.
\end{equation}
От направления $s$ он зависит только через $s^2 = 1$, а от скорости $k$ не зависит вовсе (множитель $k^2 s^2$ сокращается): $q$ --- функция одной фазы. Ключевое свойство: $q(\mu)$ \textbf{меняет знак}, с кроссовером при $\mu_{\text{cross}} \approx 0.874$:

\begin{table}[htbp]
\centering
\small
\begin{tabular}{@{}lccc@{}}
\toprule
\textbf{состояние} & $\mu$ & $q$ & $|H_\tau|$ (ед.\ $k$) \\
\midrule
каша & $\to 1$ & $+1$ & $\to \infty$ \\
переход & $0.874$ & $0$ & $\approx 5.2$ \\
сегодня & $\approx 0.76$ & $-0.43$ & $\approx 3.4$ \\
середина & $0.5$ & $-0.84$ & $2.33$ \\
отскок объектов & $\to 0$ & $-1$ (де Ситтер) & $2$ \\
\bottomrule
\end{tabular}
\end{table}

Переход замедление $\to$ ускорение возникает из самой кривизны $b(\mu)$, без введённой постоянной, и его механизм читается из двухполюсного разложения редшифта (\S5.2). Вблизи каши ($\mu_e \to 1$) доминирует пространственный множитель: расширение определяется растяжением пространства, проступающего сквозь монолитную объектность, --- режим замедления ($q \to +1$). В поздней Вселенной ($\mu_e \to \mu_0$) доминирует объектная шкала наблюдателя --- режим ускорения ($q_0 \approx -0.43$ при $\mu_0 \approx 0.76$), с асимптотикой де Ситтера $q \to -1$ при $\mu \to 0$. Ускорение здесь --- не результат того, что что-то «толкает» Вселенную: объектная линейка наблюдателя и масштаб пространства расходятся всё быстрее, и их относительный разбег считывается изнутри как ускоренное расширение. Ускорение --- тень разбега, а не сила; де Ситтер --- асимптотика приближения к отскоку объектов, а не вводимая постоянная.

Часовой множитель $d\tau = dt/\mu$ --- не косметика: он существенно меняет кинематику по всему диапазону, а не только у краёв. В середине диапазона ($\mu = 0.5$) параметр замедления без учёта часов равен $\approx -1.27$ против $\approx -0.84$ с часами.\footnote{Формула без часового множителя (производные прямо по $dt$, где $d\mu/dt = s\cdot k$): $q_{\text{noclock}}(\mu) = (-2\mu^4 + 9\mu^2 - 6)/(\mu^4 - 4\mu^2 + 4)$; при $\mu = 0.5$ она даёт $-1.27$ против $-0.84$ по полной формуле \S5.4.}

\subsection{Расстояния}

Свет идёт с $c = \text{const}$ в локальных единицах каждой эпохи; сопутствующее расстояние накапливается вдоль пути как
\begin{equation}
d\chi = \frac{c}{k} \cdot \frac{b(\mu_0)}{b(\mu)} \cdot \frac{d\mu}{\mu}\,,
\end{equation}
интегрируется от фазы излучения $\mu_e$ вниз до фазы приёма $\mu_0$ (вдоль пути света $\mu$ убывает от $\mu_e$ к $\mu_0$). С $b(\mu) = \sqrt{1 - \mu^2}/\mu^2$ интеграл берётся в замкнутом виде ($\int \mu/\sqrt{1 - \mu^2}\, d\mu = -\sqrt{1 - \mu^2}$):
\begin{align}
\chi(\mu_e; \mu_0) &= \frac{c}{k} \cdot \frac{\sqrt{1 - \mu_0^2}}{\mu_0^2} \cdot \left[ \sqrt{1 - \mu_0^2} - \sqrt{1 - \mu_e^2} \right]\,, \\
d_L(z) &= (1 + z)\cdot\chi(\mu_e; \mu_0)\,, \qquad 1 + z = \frac{b(\mu_0)}{b(\mu_e)}\,, \\
\text{DM}(z) &= 5\log_{10}(d_L / 10\,\text{пк})\,.
\end{align}
Два следствия. \textbf{Горизонт частиц конечен сам по себе}, без отсечки: у каши
\begin{equation}
\chi(\mu_e \to 1) = \frac{c}{k}\cdot\frac{1 - \mu_0^2}{\mu_0^2} \quad \text{--- конечно.}
\end{equation}
\textbf{Самосогласованность:} тождество $D_H = d\chi/dz$ (где $D_H = c/H$) выполняется точно --- темп $H(z)$ из \S5.4 и расстояния $\chi(z)$ связаны одной геометрией, и расстояния, вычисленные обоими путями, совпадают.

\subsection{Перевод на язык метрики Фридмана}

Ядро DSIC можно прочитать и на стандартном языке общей теории относительности --- не меняя ни одной формулы, а лишь сопоставив величины. Наблюдаемый масштабный фактор $b(\mu)$ играет роль масштабного фактора $a$ в плоской метрике Фридмана--Робертсона--Уокера, где собственное время наблюдателя $d\tau = dt/\mu$ служит космическим временем:
\begin{equation}
ds^2 = -c^2\, d\tau^2 + b(\mu)^2 \cdot dx^2\,.
\end{equation}
Это тождественный перевод: $b(\mu)$ --- не аналог, а буквально масштабный фактор такой метрики. Тогда вопрос «какую эффективную среду имитирует геометрия DSIC» решается стандартным тождеством плоской FRW-космологии, связывающим уравнение состояния $w_{\text{eff}} = p/\rho$ с параметром замедления,
\begin{equation}
w_{\text{eff}} = \frac{2q - 1}{3}\,,
\end{equation}
и, подставляя $q(\mu)$ из \S5.4, даёт замкнутое выражение
\begin{equation}
w_{\text{eff}}(\mu) = \frac{-\mu^4 + (16/3)\mu^2 - 4}{\mu^4 - 4\mu^2 + 4}\,.
\end{equation}
Ни одной новой сущности здесь не введено: $w_{\text{eff}}$ целиком определяется тем же $q(\mu)$, что уже выведено из ядра. Предельные значения показывают, что DSIC на языке ОТО --- это плоская Вселенная с единственной средой, чьё эффективное уравнение состояния плавно скользит от излучения к космологической постоянной:

\begin{table}[htbp]
\centering
\small
\begin{tabular}{@{}lccl@{}}
\toprule
\textbf{состояние} & $\mu$ & $w_{\text{eff}}$ & \textbf{интерпретация в терминах ОТО} \\
\midrule
каша & $\to 1$ & $+1/3$ & как излучение (замедление) \\
переход & $0.874$ & $-1/3$ & граница ускорения ($q = 0$) \\
сегодня & $\approx 0.76$ & $-0.62$ & как тёмная энергия (не константа) \\
середина & $0.5$ & $-0.89$ & близко к де Ситтеру \\
отскок объектов & $\to 0$ & $-1$ & точно космологическая постоянная \\
\bottomrule
\end{tabular}
\end{table}

Подчеркнём смысл этого перевода. DSIC \textbf{не постулирует} среду с таким $w(z)$ --- она получает эту кривую из чистой геометрии перетекания масштабов; $w_{\text{eff}}(\mu)$ есть лишь её тень на языке ОТО. Там, где $\Lambda$CDM вводит $\Omega_\Lambda$ руками, DSIC воспроизводит эквивалентную динамику одним геометрическим параметром $\mu_0$. Перевод не отнимает у модели её онтологию --- он лишь делает её сопоставимой со стандартным формализмом и показывает, что DSIC не конфликтует с общей теорией относительности на уровне фоновой метрики, а укладывается в неё как вселенная с динамическим $w(z)$.

\begin{quote}
\textit{Замечание.} Совпадение фоновой кинематики с некоторой FRW-моделью не означает совпадения на уровне возмущений и роста структур: динамический $w(z)$ фиксирует фон, но не однозначно задаёт эволюцию неоднородностей (см.\ \S12). Перевод касается именно фона.
\end{quote}
\section{Сверка с наблюдениями ($z < 2.3$)}

Форма всех наблюдаемых зависимостей жёстко фиксирована функцией $b(\mu) = \sqrt{1-\mu^2}/\mu^2$; аналоги $\Omega_m, \Omega_\Lambda$ отсутствуют. Свободны начальное условие $\mu_0$ (фаза «сегодня») и шкала $H_0$ (для сверхновых --- ещё абсолютная величина $M$, как и в $\Lambda$CDM).

О шкалах. BAO, космические хронометры и возраст считаются при $\mu_0 = 0.76$ на ранней шкале $H_0 = 67.4$; перевод геометрии в мегапарсеки идёт через $c/k = (c/H_0)\cdot(2-\mu_0^2)/(1-\mu_0^2) \approx 14978$ Мпк на той же шкале. Прямая подгонка к сверхновым (\S6.1) даёт $H_0 \approx 73$ --- это известное «напряжение Хаббла» [8], воспроизводящееся у DSIC и $\Lambda$CDM одинаково; фитовое значение в раннюю шкалу не подставляется, поскольку это разные шкалы, разведённые самим напряжением. Модель не создаёт нового напряжения и не устраняет существующего.

\subsection{Сверхновые: Pantheon+SH0ES}

Модель сверена с полным каталогом \textbf{Pantheon+SH0ES} (1701 SN; в фит идут 1657, включая 77 калибраторов [5] как якорь абсолютной шкалы) с полной ковариационной матрицей STAT+SYS. Свободны три параметра ($\mu_0, H_0, M$); плоская $\Lambda$CDM фитится на тех же данных с тремя же ($\Omega_m, H_0, M$):
\begin{align*}
\text{DSIC}: &\quad \mu_0 = 0.7600 \pm 0.0091\,, \;\; H_0 = 73.35 \pm 1.01\,, \;\; M = -19.243 \pm 0.030\,, \;\; \chi^2/\text{dof} = 0.8783 \\
\Lambda\text{CDM}: &\quad \Omega_m = 0.3308 \pm 0.0180\,, \;\; H_0 = 73.53 \pm 1.02\,, \;\; M = -19.244 \pm 0.030\,, \;\; \chi^2/\text{dof} = 0.8783
\end{align*}
\begin{equation*}
\Delta\chi^2\;(\text{DSIC} - \Lambda\text{CDM}) = +0.05 \qquad (\Delta\text{AIC} = \Delta\text{BIC} = +0.05)
\end{equation*}
Погрешности --- $1\sigma$ из гессиана $\chi^2$ в минимуме ($C = 2H^{-1}$); параметр модели измерен с точностью ${\sim}1.2\%$. Отдельная валидация конвейера: $\Lambda$CDM-ветка того же пайплайна воспроизводит опубликованные значения анализа Pantheon+SH0ES ($\Omega_m \approx 0.33 \pm 0.018$, $H_0 \approx 73.5 \pm 1.0$) --- обе модели считаются одним кодом на одних данных, и совпадение контрольной ветки с литературой подтверждает корректность сборки.

Разница качества подгонки --- пять сотых: на полностью откалиброванных данных DSIC \textbf{статистически неотличима} от $\Lambda$CDM ($|\Delta\chi^2| \ll 2$). Геометрия одного параметра $\mu_0$ воспроизводит ту же зависимость расстояния от красного смещения, которую стандартная модель задаёт двумя параметрами плотности, --- без тёмной энергии. Постоянная Хаббла при этом выходит как настоящий вывод из якоря ($H_0 \approx 73$), а не как поглотитель сдвига $M$. Прогон со STAT-ONLY ковариацией даёт $\mu_0 = 0.7661$, $\Delta\chi^2 = -0.39$ --- результат устойчив к способу учёта ошибок.

\subsection{Барионный масштаб: DESI DR2}

Данные DESI DR2 (Table IV [6], $0.5 < z < 2.33$) дают раздельно $D_M/r_d$ и $D_H/r_d$ в шести бинах. Безразмерное отношение $D_M/D_H$ не зависит ни от $r_d$, ни от $H_0$ и тестирует чистую геометрию расширения:
\begin{equation*}
D_M/D_H \;(\text{DSIC vs DESI DR2, 6 точек}): \quad \chi^2/\text{dof} \approx 0.78
\end{equation*}
--- наравне со стандартной моделью на тех же точках. Точка $z = 0.934$ даёт наибольшее напряжение и в $\Lambda$CDM (${\approx}\,2.5\sigma$) --- это особенность данных, а не модели.

Масштаб $r_d$ восстановлен обратной задачей: $r_d = D_M^{\text{DSIC}}/(D_M/r_d)_{\text{DESI}}$ и независимо $r_d = D_H^{\text{DSIC}}/(D_H/r_d)_{\text{DESI}}$ в каждом бине. Результат постоянен по $z$ и согласован между поперечным и радиальным измерениями:
\begin{equation*}
r_d \approx 148\;\text{Мпк} \quad (\text{разброс по } z \approx 1.4\%;\; \text{по } D_M \approx 147.7,\; \text{по } D_H \approx 150.2)\,,
\end{equation*}
что совпадает со стандартным масштабом ${\approx}\,147$ Мпк. Динамической зависимости $r_d(z)$ в данных нет: полюса $P$ и $Q$ меняются с $z$, но восстановленный физический масштаб остаётся постоянным --- линейка DSIC статична в наблюдаемом окне.

\subsection{Темп расширения и возраст}

Космические хронометры [9] $H(z)$ на $0.07 < z < 2$ дают $\chi^2/\text{dof} \approx 0.37$, наравне со стандартной моделью. Интеграл по истории расширения даёт $t_0\cdot H_0 \approx 0.925$, то есть при $H_0 = 67.4$ возраст $t_0 \approx 13.4$ Гыр. Оговорка о сравнении: часто цитируемые $13.80 \pm 0.02$ Гыр --- возраст, \textit{выведенный внутри $\Lambda$CDM} из данных Planck, и мерить им не-$\Lambda$CDM модель отчасти циклично; модельно-независимые звёздные возрасты (старейшие шаровые скопления, ${\sim}13.5 \pm 0.5$ Гыр) значение $t_0 \approx 13.4$ Гыр проходят.

\subsection{Параметр замедления и красное смещение перехода}

Через связь $\mu(z)$ (\S5.2) при $\mu_0 \approx 0.76$ теоретическая кривая $q(z)$ сопоставлена с модельно-независимыми кинематическими реконструкциями [12]:

\begin{table}[htbp]
\centering
\small
\begin{tabular}{@{}ccc@{}}
\toprule
$z$ & $q_{\text{DSIC}}$ & \textbf{реконструкция (модельно-незав.)} \\
\midrule
$0$ & $-0.43$ & $-0.55$ \\
$0.3$ & $-0.25$ & $-0.34$ \\
$0.65$ & $-0.06$ & $0.0$ \\
$1.0$ & $+0.11$ & $+0.18$ \\
$1.5$ & $+0.30$ & $+0.33$ \\
$2.0$ & $+0.44$ & $+0.42$ \\
\bottomrule
\end{tabular}
\end{table}

Форма $q(z)$ --- S-образная, с замедлением на высоком $z$ и ускорением на низком --- воспроизводит реконструкции в пределах разброса; красное смещение перехода $z_t \approx 0.77$ [10] согласуется с наблюдаемым $z_t \approx 0.6\text{--}0.8$.

\subsection{Совместный фит и $Om(z)$}

Совместный фит SN+BAO формально слабо предпочитает DSIC ($\Delta\chi^2 \approx -5.0$); преимущество частично отражает известное натяжение DESI--$\Lambda$CDM и не интерпретируется как решающее. $Om(z)$-диагностика на текущих данных моделей \textbf{не разрешает}: знак $\Delta\chi^2$ зависит от нормировки $H_0$. Это null-результат.

\subsection{Итог сверки}

Единое начальное условие $\mu_0 \approx 0.76$ согласованно описывает шесть независимых проб поздней Вселенной:

\begin{table}[htbp]
\centering
\small
\begin{tabular}{@{}p{2.6cm}p{4.0cm}p{5.2cm}@{}}
\toprule
\textbf{проба} & \textbf{данные} & \textbf{результат} \\
\midrule
сверхновые, $d_L(z)$ & Pantheon+SH0ES (1657 SN, якорь + ковариация) & $\Delta\chi^2 = +0.05$ vs $\Lambda$CDM --- неотличимо \\
\addlinespace
BAO $D_M/D_H$ & DESI DR2, 6 точек & $\chi^2/\text{dof} \approx 0.78$ \\
\addlinespace
линейка $r_d$ & DESI DR2 $D_M/r_d$, $D_H/r_d$ & ${\approx}\,148$ Мпк, статична \\
\addlinespace
хронометры $H(z)$ & $z < 2$ & $\chi^2/\text{dof} \approx 0.37$ \\
\addlinespace
возраст $t_0$ & интеграл по истории & ${\approx}\,13.4$ Гыр (звёздные возрасты ${\sim}13.5 \pm 0.5$) \\
\addlinespace
переход $z_t$ & реконструкции $q(z)$ & ${\approx}\,0.77$ (набл.\ $0.6\text{--}0.8$) \\
\bottomrule
\end{tabular}
\end{table}

Положение наблюдателя на цикле фиксируется этими наблюдениями согласованно: одна геометрия, один параметр, без тёмной энергии.

\section{Отличимые предсказания и фальсифицируемость}

На поздних данных DSIC и $\Lambda$CDM практически вырождены; разводит их высокий $z$, где модель делает резкие проверяемые предсказания:
\begin{itemize}
  \item \textbf{форма $Om(z)$} на расширенном диапазоне красных смещений;
  \item \textbf{радиальный BAO в области Lyman-$\alpha$} ($z \approx 3.5$): расхождение с $\Lambda$CDM до ${\sim}5\%$;
  \item \textbf{темп расширения $E(z)$ при $z = 5\text{--}10$}: расхождение $+15\ldots+47\%$.
\end{itemize}
Важно разграничить источники сигналов: расхождение по $D_H$ в окне $2.33 < z < 4.53$ --- след \textbf{ядра} ($D_H = c/H$ зависит только от $\mu_0$) и от механики Части II не зависит; след порога $\mu_d$ лежит в $D_M$ выше $z \approx 4.5$ и обсуждается в \S11.5.

\section{Границы применимости и открытые направления}

Заявления намеренно держатся узко: сделанным считается только фоновая кинематика поздней Вселенной ($z < 2.3$) из \S6. Не сделано следующее.

\medskip
\noindent\textbf{Рост структур.} Возмущения выведены не из ядра, но фон DSIC проверен на совместимость со стандартным (GR-предельным) ростом линейных неоднородностей: на компиляции $f\sigma_8$ модель описывает данные не хуже $\Lambda$CDM (\S12). Полноценный вывод уравнения роста из связанности (Приложение B) остаётся открытым.

\medskip
\noindent\textbf{Ранняя Вселенная.} Из ядра расстояние до реликта недомеряет ${\sim}21\%$. Это закрывается вторым этажом модели (Часть II): форма поправки выводится из постулата о связанности объектов (P6, \S10), а её единственный порог $\mu_d$ калибруется по одному наблюдению. Что это решает, что нет и почему --- \S11.6.

\medskip
\noindent\textbf{Онтология против анзаца.} Вид $b(\mu) = \sqrt{1-\mu^2}/\mu^2$ выводится здесь из круговой геометрии полюсов, однако читатель вправе принять его и просто как удачный однопараметрический анзац: формулы работают сами по себе, и онтологическая картина для их работы не обязательна.

\medskip
\noindent\textbf{Гравитация.} Локальное распределение проступающего пространства требует той же недостающей ядру сущности --- связанности объектов --- и разобрано в Приложении B как позднее проявление второго этажа; сверки с данными у этого модуля пока нет.
\section{Соотношение с $\Lambda$CDM}

DSIC не «опровергает» $\Lambda$CDM. На поздних данных модели статистически неотличимы; DSIC предлагает другую \textbf{онтологию} того же фона --- геометрию вместо тёмной энергии --- при равном числе свободных параметров. Ценность не в качестве подгонки (оно не лучше), а в происхождении кривой: та же зависимость расстояния от красного смещения получается из одного геометрического начала, без $\Omega_\Lambda$, причём переход замедление $\to$ ускорение и конечность горизонта частиц выводятся, а не постулируются. Решающая проверка --- измерения высокого $z$ из \S7.

\bigskip
\hrule
\bigskip

\part{Второй этаж модели: связанность объектов}

\begin{quote}
\textbf{Статус.} Всё в этой части опирается на ядро, но не выводится из него: в полюсах $S, O$ есть только перетекание масштаба, а обеим конструкциям ниже нужна вторая, независимая сущность --- \textbf{связанность объектов}. Ядро (Часть I) самодостаточно: все шесть проб \S6 от Части II не зависят, и поправка второго этажа тождественно равна нулю во всей области их данных. Часть II вводит одну новую сущность, один дополнительный постулат и одну измеренную константу; где заканчивается вывод и начинается калибровка --- оговаривается на каждом шаге.
\end{quote}

\section{Вторая сущность и постулат отлипания}

К одной и той же недостающей сущности ядро упирается с двух сторон. Со стороны ранней Вселенной: геометрия ядра недомеряет расстояние до реликта на ${\sim}21\%$ (\S11.1) --- не хватает физики фазы, в которой материя ещё слиплась, а «слипание» есть состояние связанности. Со стороны поздней: качественная картина, по которой пространство прибывает преимущественно между слабо связанными гравитацией объектами (Приложение B), требует меры этой связи. Обе надстройки --- раннее отлипание и позднее распределение пространства --- суть два проявления одной новой сущности: \textbf{связанности}. Вводим её один раз.

\medskip
\noindent\textbf{P6 (постулат связанности и отлипания).}
\begin{itemize}
  \item Ранняя эпоха --- почти монолитная объектность, и порог относится к \textbf{средней среде вдоль пути света}, а не к каждой точке. Пока доминирующий монолит в среднем не расступился ($\mu > \mu_d$), измерительная цепочка внутреннего наблюдателя вдоль пути не определена: линейка $L \propto \mu$ собрана из раздельных объектов, а средой всё ещё правит слипшаяся компонента. Свету в этой фазе почти негде распространяться --- путь считается не по метрике наблюдателя, а по чистой доступности пространства: добавка к пути $\propto d\mu/P = d\mu/\sqrt{1-\mu^2}$. Взвешивание $b(\mu_0)/b(\mu)$, которым взвешен путь поздней формулы (\S5.5), здесь не применяется --- взвешивать нечем.
  \item Локальные сгустки, в которых связанность сильнее всего, \textbf{падают в себя} (\S3.4) и коллапсируют в объекты задолго до порога --- первые галактики и квазары. Их появление не меняет средней доступности пространства вдоль пути света и не сдвигает порог: $\mu_d$ фиксирует момент, когда расступился монолит в среднем, а не момент рождения первых объектов.
  \item Отлипание среды --- механическое событие, а не процесс: \textbf{жёсткий порог} $\mu_d$. Когда материя впустила пространство до степени, при которой смогла отлипнуть от самой себя и расступиться, свет получает прибывающую возможность распространяться: при $\mu \le \mu_d$ поправка тождественно равна нулю, работает поздняя геометрия ядра без изменений.
  \item Форма зафиксирована принципом: амплитуда естественная ($A \equiv 1$, форма $1/P$ без множителя), окно резкое (без сглаживания). Свободным остаётся ровно одно число --- порог $\mu_d$.
\end{itemize}

Разграничение «средняя среда vs первые островки» --- не косметика, а необходимое условие согласованности: галактики JWST наблюдаются до $z \approx 10\text{--}14$ ($\mu_e(z{=}14) \approx 0.9972$), то есть \textbf{выше} порога $\mu_d = 0.9806$ ($z_{\text{detach}} \approx 4.53$). Буквальное чтение «раздельных объектов не существует при $\mu > \mu_d$» противоречило бы этим наблюдениям; в принятой формулировке первые объекты --- сгустки, упавшие в себя внутри монолита там, где связанность максимальна, и их существование с порогом средней среды совместимо. Проверочный скрипт выводит этот флаг согласованности явно.

Физическая картина --- аналогия пресса. Набор предметов сжали; когда пресс отпускает, предметы сначала ещё распирают друг друга, а в какой-то момент \textbf{отлипают} --- и после этого давления на стенки больше нет, хотя пресс продолжает расходиться. Точка отлипания и есть $\mu_d$ --- момент, когда расступается масса в целом; отдельные предметы в глубине могли сцепиться между собой раньше, на порог это не влияет. Резкость порога мотивирована именно механикой: отлипание --- событие, гладкое окно вводило бы лишний параметр без физического основания.

Деталь в поддержку естественности формы. Базовый интегранд пути ядра равен $(\sqrt{1-\mu_0^2}/\mu_0^2)\cdot\mu/P$ (\S5.5), интегранд поправки --- $1/P$; у самой каши ($\mu \to 1$) они совпадают по форме: оба $\propto 1/P$. Добор в слипшейся фазе асимптотически сопоставим с базовой мерой пути --- свет «пробирается» и сквозь объектность, а не только по уже проступившему пространству.

\medskip
\noindent\textbf{Статус $\mu_d$.} Значение порога из ядра не выводится: в полюсах $S, O$ нет числа, выделяющего $0.9806$, --- порог есть условие на связанность, которой в ядре нет. Но при форме, зафиксированной P6, вырождение «амплитуда $\leftrightarrow$ порог» ломается принципом, и $\mu_d$ становится \textbf{измеренной константой второго этажа} --- того же статуса, что $\mu_0$ в ядре. Теории разрешено иметь измеренные константы; не разрешено --- произвольные функции. Без P6 калибровка была бы вырождена, что честно показано таблицей в \S11.2.
\section{Ранняя Вселенная: порог отлипания}

\subsection{Недобор ядра}

Параметры ядра на ранней шкале --- те же, что в \S6 для BAO, хронометров и возраста: $\mu_0 = 0.76$, $H_0 = 67.4$, $c/k = 14978.2$ Мпк. Опорная цифра Planck 2018 [7] (TT,TE,EE+lowE+lensing): $z^* = 1089.92$, $r^* = 144.43$ Мпк, $100\theta^* = 1.04110$, откуда сопутствующее расстояние до поверхности последнего рассеяния $d_{\text{LSS}} = r^*/\theta^* = 13872.8 \pm 25$ Мпк (внутренняя согласованность $r^*/\theta^*$ сверяется самопроверкой скрипта). Геометрия ядра без поправки:
\begin{align*}
\mu_{\text{cmb}}(z^* = 1089.92) &= 0.9999994681 \qquad (\text{из } 1+z = b(\mu_0)/b(\mu_e)) \\
D_M \text{ база до CMB} &= 10936.2\;\text{Мпк} \\
d_{\text{LSS}}\;(\text{Planck}) &= 13872.8\;\text{Мпк} \\
\text{недобор} &= 2936.6\;\text{Мпк} \;\;(21.2\%)
\end{align*}
Причина недобора --- сторона ядра: горизонт частиц конечен «бесплатно» (\S5.5), почти весь путь набирается уже к $z \approx 10$, дальше интеграл расстояния выдыхается. Ближняя и средняя Вселенная ($z < 2.3$) при этом считается верно, поэтому порог отлипания воздействует только на дальний конец.

\subsection{Поправка и калибровка порога}

По P6 добавка к пути в слипшейся фазе есть $1/P$, отсечённая порогом; её интеграл --- $\arcsin\mu$:
\begin{align*}
\text{поправка}(\mu) &= \frac{1}{\sqrt{1-\mu^2}} \;\; \text{при } \mu > \mu_d\,; \qquad 0 \;\; \text{при } \mu \le \mu_d \\[4pt]
D_M^{\text{испр}}(\mu_e) &= D_M(\mu_e) + (c/k)\cdot[\arcsin\mu_e - \arcsin\mu_d] \quad \text{при } \mu_e > \mu_d \\
D_M^{\text{испр}}(\mu_e) &= D_M(\mu_e) \hspace{4.9cm} \text{при } \mu_e \le \mu_d
\end{align*}
Порог фиксируется требованием точного замыкания недобора:
\begin{align*}
&(c/k)\cdot[\arcsin\mu_{\text{cmb}} - \arcsin\mu_d] = 2936.6\;\text{Мпк} \\
&\Rightarrow\; \mu_d = 0.980640\,, \qquad z_{\text{detach}} = 4.526 \\
&D_M^{\text{испр}}(\mu_{\text{cmb}}) = 13872.8\;\text{Мпк} \quad (= d_{\text{LSS}}\text{ Planck, замыкание точное})
\end{align*}
Для честности покажем, что решает P6. Без фиксации амплитуды уравнение $A\cdot(c/k)\cdot[\arcsin\mu_{\text{cmb}} - \arcsin\mu_d] = \text{недобор}$ имеет кривую решений --- любая амплитуда даёт своё $\mu_d$ с тем же расстоянием до CMB:

\begin{table}[htbp]
\centering
\small
\begin{tabular}{@{}ccc@{}}
\toprule
$A$ & $\mu_d$ & $z_{\text{detach}}$ \\
\midrule
0.5 & 0.923706 & 1.506 \\
0.8 & 0.969868 & 3.344 \\
\textbf{1.0} & \textbf{0.980640} & \textbf{4.526} \; $\leftarrow$ фиксировано P6 \\
1.2 & 0.986514 & 5.691 \\
2.0 & 0.995097 & 10.266 \\
5.0 & 0.999190 & 26.922 \\
\bottomrule
\end{tabular}
\end{table}

Одно наблюдение фиксирует один параметр: степеней свободы ноль. Значение $z_{\text{detach}} \approx 4.53$ --- калибровка при принципиально зафиксированной форме, а не предсказание; предсказанием оно не предъявляется.

\subsection{Свойства оптической поправки}

Поправка --- \textbf{оптическая}, а не динамическая: удлиняется путь света, темп расширения $H(z)$ не меняется. Из этого выбора следуют проверяемые свойства, и все они здесь выполняются по построению:
\begin{itemize}
  \item \textbf{Частота не затрагивается.} Красное смещение остаётся $1+z = b(\mu_0)/b(\mu_e)$ (\S5.2); удлинение пути на него не влияет.
  \item \textbf{Дуальность расстояний сохраняется.} $d_L$ и $d_A$ строятся из одной и той же $\chi_{\text{испр}}$, поэтому соотношение Этерингтона $d_L = (1+z)^2\cdot d_A$ выполняется автоматически.
  \item \textbf{Ахроматичность и недиссипативность --- жёсткое требование.} Чернотельность спектра реликта (FIRAS) исключает любую «среду», зависящую от частоты или рассеивающую; поправка обязана быть геометрией пути, а не взаимодействием, --- форма $1/P$ этому удовлетворяет.
  \item \textbf{Горизонт частиц остаётся конечным.} С поправкой $\chi_{\text{испр}}(\mu_e \to 1) = (c/k)\cdot[(1-\mu_0^2)/\mu_0^2 + (\pi/2 - \arcsin\mu_d)] \approx 13906$ Мпк --- всего на ${\sim}33$ Мпк дальше поверхности последнего рассеяния: в исправленной геометрии реликт лежит почти вплотную к горизонту частиц.
\end{itemize}

\subsection{Поздняя Вселенная не затронута}

Поправка строго равна нулю ниже порога --- это тождество, а не приближение:

\begin{table}[htbp]
\centering
\small
\begin{tabular}{@{}ccc@{}}
\toprule
$z$ & $D_M$ база, Мпк & поправка, Мпк \\
\midrule
0.50 & 1932.8 & $+0.0$ \\
1.00 & 3384.2 & $+0.0$ \\
2.33 & 5792.7 & $+0.0$ \\
4.00 & 7335.6 & $+0.0$ \\
4.53 & 7653.3 & $+0.0$ \\
4.63 & 7707.8 & $+49.3$ $\leftarrow$ включена \\
5.00 & 7896.9 & $+220.4$ $\leftarrow$ включена \\
10.00 & 9247.2 & $+1433.0$ $\leftarrow$ включена \\
\bottomrule
\end{tabular}
\end{table}

Все шесть бинов DESI DR2 ($z \le 2.33$) поправки не чувствуют, $\chi^2/\text{dof}$ на BAO не меняется --- результаты \S6 сохранены полностью.

\subsection{След порога и досягаемость зондов}

Поправка к $D_M$ становится заметной только глубоко за порогом:

\begin{table}[htbp]
\centering
\small
\begin{tabular}{@{}ccc@{}}
\toprule
$z$ & объект & поправка к $D_M$ \\
\midrule
4.53 & порог отлипания & $+0.0\%$ \\
5.00 & галактики JWST & $+2.8\%$ \\
6.00 & квазары & $+7.1\%$ \\
8.00 & реионизация & $+12.4\%$ \\
11.00 & первые галактики & $+16.6\%$ \\
1089.92 & CMB (калибровка) & $+26.9\%$ \\
\bottomrule
\end{tabular}
\end{table}

Досягаемость зондов: DESI DR2 ($z \le 2.33$), DESI DR3 ($\sim$2027, $z \le 3.55$) и DESI-II/LBG ($\sim$2029, $2.5 < z < 3.5$) --- все ниже порога. Расхождение $D_H/r_d$ в окне $2.33 < z < 4.53$ (до ${\sim}5\%$ на $z \approx 3.5$, \S7) --- след ядра, а не $\mu_d$. Сам порог влияет лишь на $D_M$ выше $z \approx 4.5$, где точных стандартных линеек, кроме самой CMB-калибровки, нет: второго независимого якоря для $\mu_d$ в досягаемых данных не существует --- до появления точной шкалы расстояний на $z \gtrsim 5\text{--}6$.

\subsection{Что второй этаж решает, а что нет}

Закрывается ровно одно число --- расстояние до поверхности последнего рассеяния, точно и без последствий для поздней Вселенной. Открытым остаётся следующее, и это принципиально.

\medskip
\noindent\textbf{Заимствованный звуковой горизонт.} Цель замыкания $d_{\text{LSS}} = r^*/\theta^*$ использует $r^* = 144.43$ Мпк --- величину, не измеренную напрямую, а вычисленную внутри $\Lambda$CDM из физики фотон-барионной плазмы, которой в DSIC нет. Честная постановка: прямо измеренной величиной является угол $\theta^*$, на него и идёт калибровка; происхождение самого акустического масштаба ${\sim}147$ Мпк --- согласованно проступающего и в поздней линейке $r_d \approx 148$ Мпк (\S6.2) --- объявляется открытым вопросом ядра.

\medskip
\noindent\textbf{Темп расширения не тронут --- и это самое узкое место, заявляемое числом.} Оптическая поправка меняет путь, но не $H(z)$: при $z = 5\text{--}10$ темп остаётся выше $\Lambda$CDM на $+15\ldots+47\%$ (\S7), а значит, космическое время на этих эпохах сжато. Количественно (общая шкала $H_0 = 67.4$; $\Omega_m = 0.303$ из совместного фита \S6, отчего $t_0$ $\Lambda$CDM здесь 13.95 Гыр --- обе модели поставлены в равные условия; $t_0$ DSIC $= 13.38$ Гыр):

\begin{table}[htbp]
\centering
\small
\begin{tabular}{@{}cccc@{}}
\toprule
$z$ & $t_{\text{DSIC}}$, Гыр & $t_{\Lambda\text{CDM}}$, Гыр & отношение \\
\midrule
3.0 & 1.73 & 2.18 & 0.79 \\
5.0 & 0.81 & 1.19 & 0.68 \\
6.0 & 0.60 & 0.95 & 0.64 \\
7.5 & 0.42 & 0.71 & 0.59 \\
10.0 & 0.25 & 0.48 & 0.52 \\
12.0 & 0.18 & 0.37 & 0.48 \\
\bottomrule
\end{tabular}
\end{table}

Солпитеровская прикидка обостряет картину. На эддингтоновском пределе e-складывание массы чёрной дыры занимает ${\sim}45$ млн лет; рост от звёздного зародыша ${\sim}10^2\,M_\odot$ до наблюдаемых у квазаров $z \approx 7.5$ масс ${\sim}10^9\,M_\odot$ требует ${\sim}16$ e-складываний $\approx 0.72$ Гыр. $\Lambda$CDM располагает 0.71 Гыр --- впритык (что уже само по себе мотивирует тяжёлые зародыши); DSIC располагает 0.42 Гыр $\approx 9$ e-складываний --- обязательны либо зародыши прямого коллапса ${\sim}10^5\,M_\odot$ при непрерывной аккреции, либо супер-эддингтоновские режимы. Это не опровержение --- супер-эддингтоновская аккреция обсуждается всерьёз, в том числе из-за «перемассивных» чёрных дыр, обнаруживаемых JWST, --- но это ближайшая реальная проверка ядра и, вероятно, первая точка серьёзного сопротивления модели; она предъявляется числом, а не оговоркой.

Картина ядра предлагает и естественный, хотя и не выводимый строго, механизм под эти наблюдения. Отлипание монолита (\S10) --- событие \textbf{средней среды}, а не каждой точки: когда пространство проступает и объектность в целом расступается, отдельные участки предельной плотности могут по разным причинам \textbf{не расступиться} и сохранить плотность каши --- как во всяком разжимающемся веществе часть сгустков остаётся спрессованной. Такие не впустившие пространство области --- готовые сверхтяжёлые зародыши чёрных дыр, существующие \textbf{с самого начала разлёта}, а не наращённые аккрецией из звёздных остатков. В этой картине «перемассивные» чёрные дыры, обнаруживаемые JWST на больших $z$, --- не аномалия, требующая рекордно быстрого роста, а ожидаемый след неоднородного отлипания: где-то монолит расступился, а где-то остался кашей. Строгого вывода распределения таких зародышей ядро не даёт (для этого нужна статистика связанности, Приложение B), но само их существование картине не противоречит, а следует из неё, и сжатый бюджет времени $z = 5\text{--}12$ перестаёт быть чистой уязвимостью: часть массивных ранних дыр модель ожидает «врождёнными».

\medskip
\noindent\textbf{Реликт --- не одно число.} Planck ограничивает всю структуру акустических пиков, хвост затухания и поляризацию --- тысячи мультиполей. Второй этаж закрывает один параметр ($\theta^*$); всё остальное требует теории возмущений и состава (барионы, фотоны, нейтрино), которых в модели нет.

\subsection{Трактовка реликтового отпечатка}

Отпечаток --- стоп-кадр формы объектности на поздней стадии каши, когда пространство уже начало проступать и свет впервые смог выйти из монолита: контраст на карте есть контраст плотности проступающего пространства. Каша при этом --- не строго нулевое состояние, а отступ от нуля: $\mu_{\text{cmb}} \approx 0.9999995$, пространство уже чуть проступило. Знак контраста в этой картине: \textbf{холодные/тёмные участки $\leftrightarrow$ плотная объектность (монолит), горячие $\leftrightarrow$ проступившие поры пространства}; малость контраста ${\sim}10^{-5}$ качественно согласуется с тем, что у порога поры пространства только зарождаются. Отметим для трезвости: этот знак совпадает со стандартным крупномасштабным эффектом Сакса--Вольфа для адиабатических возмущений (сверхплотные области выглядят холоднее) --- согласованность приятная, но не свидетельство, поскольку амплитуда и спектр из модели не выведены. Наконец, важно различать: $\mu_{\text{cmb}}$ --- наблюдаемое «дно» ранней зоны, $\mu_d$ ($z \approx 4.53$) --- её поздний край, порог отлипания.
\section{Рост структур: проверка на фоне DSIC}

Ранее (\S8) рост структур был заявлен как не входящий в модель. Здесь делается минимальный, честно ограниченный шаг: проверяется, совместим ли \textbf{фон} DSIC со стандартным ростом линейных возмущений --- не выводя уравнение роста из ядра, а беря его в GR-пределе на геометрии DSIC. Это проверка на согласованность, а не вывод.

На фоне с темпом $H(z)$ (\S5.4) плотностный контраст материи $\delta$ в суб-горизонтном GR-пределе подчиняется стандартному уравнению
\begin{equation}
\frac{d^2\delta}{d(\ln a)^2} + \left[2 + \frac{d\ln H}{d\ln a}\right]\frac{d\delta}{d(\ln a)} - \frac{3}{2}\,\Omega_m(a)\,\delta = 0\,,
\end{equation}
где $\Omega_m(a) = \Omega_{m0}\cdot a^{-3}/E(a)^2$, $E = H/H_0$. Поскольку ядро плотностей не вводит, $\Omega_{m0}$ входит здесь как параметр роста. Из решения получаем скорость роста $f = d\ln\delta/d\ln a$ и наблюдаемую комбинацию $f\sigma_8(z) = f(z)\cdot\sigma_{8,0}\cdot\delta(z)/\delta(0)$, сравниваемую с измерениями искажений в пространстве красных смещений (RSD).

Сверка выполнена на компиляции \textbf{«Gold-2017»} (18 слабо коррелированных точек $f\sigma_8$; для трёх точек WiggleZ [15] при $z = 0.44, 0.60, 0.73$ учтена их ковариационная матрица). Ключевая особенность данных $f\sigma_8$ --- вырождение: они ограничивают не $\Omega_{m0}$ и $\sigma_8$ по отдельности, а их комбинацию $S_8 = \sigma_8\cdot\sqrt{\Omega_{m0}/0.3}$. Поэтому результат докладывается через $S_8$:
\begin{align*}
&\text{DSIC (GR-предел на фоне DSIC), } \mu_0 = 0.76: \\
&\quad S_8 = 0.761 \pm 0.013 \qquad (\text{устойчиво вдоль всей долины вырождения}) \\
&\quad \text{при фиксированном } \Omega_{m0} = 0.30:\;\; \sigma_8 = 0.770,\;\; \chi^2/\text{dof} \approx 0.73
\end{align*}
Прямое сравнение с $\Lambda$CDM на тех же 18 точках даёт $S_8 = 0.760$, $\chi^2/\text{dof} \approx 0.74$; разность $\Delta\chi^2 \approx +0.1$ --- статистически неотличимо. Форма $f\sigma_8(z)$ DSIC пологая, с максимумом около $z \approx 0.5$, и на текущих данных описывает измерения не хуже стандартной модели. Это не доказывает, что рост в DSIC обязан идти по GR-предельному уравнению --- ядро его не диктует, --- но снимает возражение о грубой несовместимости: фоновая геометрия DSIC совместима с наблюдаемым ростом структур при разумных $\Omega_{m0}, \sigma_8$, причём восстановленное $S_8 \approx 0.76$ согласуется со стандартным значением.

Оговорка о систематике: поправка Alcock--Paczynski за различие фоновых расстояний DSIC и опорной космологии каждого RSD-измерения не применялась; при $z < 1$ этот фактор обычно составляет несколько процентов и мал против ошибок данных. Строгий вывод уравнения роста из связанности (Приложение B) остаётся направлением развития.

\section{Заключение}

Из четырёх структурных элементов --- относительности масштаба, сохранения нормы двух полюсов, единственного равномерного потока и модульного доступа внутреннего наблюдателя --- однозначно следует масштабный фактор $b(\mu) = \sqrt{1-\mu^2}/\mu^2$, а из него вся фоновая кинематика: красное смещение $1+z = b(\mu_0)/b(\mu_e)$, темп $|H_\tau| = (2-\mu^2)/(1-\mu^2)$, меняющий знак параметр замедления $q(\mu)$ с кроссовером $\mu_{\text{cross}} \approx 0.874$, замкнутые расстояния и конечный горизонт частиц. Единственный параметр $\mu_0 \approx 0.76$ согласованно проходит шесть независимых проб поздней Вселенной, оставаясь статистически неотличимым от $\Lambda$CDM на её лучших данных и резко отличимым от неё на высоком $z$. Ускорение расширения в этой картине --- не действие скрытой компоненты, а тень разбега двух масштабов, считываемая наблюдателем, который сам сделан из одного из них.

Второй этаж добавляет к этому одну сущность --- связанность объектов --- и одну измеренную константу. Постулат отлипания P6 фиксирует форму поправки к пути света принципом (свету негде распространяться, пока средняя среда не расступилась; локальные сгустки падают в себя раньше), калибровка по реликту фиксирует порог $\mu_d \approx 0.9806$; поздняя Вселенная при этом не затрагивается тождественно, расстояние до поверхности последнего рассеяния замыкается точно, а дуальность расстояний и ахроматичность выполняются по построению. Параметр ядра измерен: $\mu_0 = 0.7600 \pm 0.0091$. Открытыми остаются происхождение акустического масштаба ${\sim}147$ Мпк, сжатое космическое время при $z = 5\text{--}12$ (при $z = 7.5$ --- 0.42 Гыр против 0.71 у $\Lambda$CDM, \S11.6) и рост структур --- именно там модель встретит решающие проверки.

\begin{thebibliography}{99}

\bibitem{riess1998} A.~G. Riess et al., \textit{Observational Evidence from Supernovae for an Accelerating Universe and a Cosmological Constant}, AJ \textbf{116}, 1009 (1998).

\bibitem{perlmutter1999} S.~Perlmutter et al., \textit{Measurements of $\Omega$ and $\Lambda$ from 42 High-Redshift Supernovae}, ApJ \textbf{517}, 565 (1999).

\bibitem{scolnic2022} D.~M. Scolnic et al., \textit{The Pantheon+ Analysis: The Full Dataset and Light-Curve Release}, ApJ \textbf{938}, 113 (2022); arXiv:2112.03863.

\bibitem{brout2022} D.~Brout et al., \textit{The Pantheon+ Analysis: Cosmological Constraints}, ApJ \textbf{938}, 110 (2022); arXiv:2202.04077.

\bibitem{riess2022} A.~G. Riess et al., \textit{A Comprehensive Measurement of the Local Value of the Hubble Constant} (SH0ES), ApJL \textbf{934}, L7 (2022); arXiv:2112.04510.

\bibitem{desi2025} DESI Collaboration (M.~Abdul-Karim et al.), \textit{DESI DR2 Results II: Measurements of Baryon Acoustic Oscillations and Cosmological Constraints}, Phys. Rev. D \textbf{112}, 083515 (2025); arXiv:2503.14738.

\bibitem{planck2018} Planck Collaboration (N.~Aghanim et al.), \textit{Planck 2018 results. VI. Cosmological parameters}, A\&A \textbf{641}, A6 (2020); erratum A\&A \textbf{652}, C4 (2021); arXiv:1807.06209.

\bibitem{divalentino2021} E.~Di Valentino et al., \textit{In the Realm of the Hubble Tension --- a Review of Solutions}, Class. Quantum Grav. \textbf{38}, 153001 (2021); arXiv:2103.01183.

\bibitem{moresco2020} M.~Moresco et al., \textit{Setting the Stage for Cosmic Chronometers} (ковариация космических хронометров), ApJ \textbf{898}, 82 (2020); arXiv:2003.07362. (С компиляцией Moresco et al. 2012; Moresco 2015; Moresco et al. 2016. Первоисточники прочих точек $H(z)$: Simon et al. 2005; Stern et al. 2010; Zhang et al. 2014; Ratsimbazafy et al. 2017; Borghi et al. 2022.)

\bibitem{farooq2013} O.~Farooq \& B.~Ratra, \textit{Hubble Parameter Measurement Constraints on the Cosmological Deceleration--Acceleration Transition Redshift}, ApJ \textbf{766}, L7 (2013); arXiv:1301.5243.

\bibitem{sahni2008} V.~Sahni, A.~Shafieloo \& A.~A. Starobinsky, \textit{Two New Diagnostics of Dark Energy} (Om-диагностика), Phys. Rev. D \textbf{78}, 103502 (2008); arXiv:0807.3548.

\bibitem{sahni2003} V.~Sahni, T.~D. Saini, A.~A. Starobinsky \& U.~Alam, \textit{Statefinder --- a New Geometrical Diagnostic of Dark Energy}, JETP Lett. \textbf{77}, 201 (2003); arXiv:astro-ph/0201498.

\bibitem{nesseris2017} S.~Nesseris, G.~Pantazis \& L.~Perivolaropoulos, \textit{Tension and Constraints on Modified Gravity Parametrizations of $G_{\text{eff}}(z)$ from Growth Rate and Planck Data}, Phys. Rev. D \textbf{96}, 023542 (2017); arXiv:1703.10538. (Компиляция «Gold-2017», 18 точек $f\sigma_8$.)

\bibitem{sagredo2018} B.~Sagredo, S.~Nesseris \& D.~Sapone, \textit{Internal Robustness of Growth Rate Data}, Phys. Rev. D \textbf{98}, 083543 (2018); arXiv:1806.10822. (Table I --- значения $z, f\sigma_8, \sigma$, использованные в \S12.)

\bibitem{blake2012} C.~Blake et al., \textit{The WiggleZ Dark Energy Survey: Joint Measurements of the Expansion and Growth History at $z < 1$}, MNRAS \textbf{425}, 405 (2012); arXiv:1204.3674. (Ковариация трёх точек $f\sigma_8$ при $z = 0.44, 0.60, 0.73$.)

\end{thebibliography}
\section*{Воспроизводимость}
\addcontentsline{toc}{section}{Воспроизводимость}

Все эмпирические утверждения статьи воспроизводятся одним самодостаточным скриптом \texttt{dsic\_test.py} из репозитория \url{https://github.com/az-zase/DSIC}

Происхождение каждой вшитой цифры задокументировано в самом файле; все опорные данные проходят самопроверки.

Блоки скрипта:

Часть I: \texttt{[0]} построчная сверка вшитого DESI DR2 с официальной Table IV; \texttt{[0b]} машинная проверка внутренних тождеств ядра (монотонность $b(\mu)$, точность $D_H = d\chi/dz$ до ${\sim}10^{-8}$, совпадение замкнутой $\chi(\mu)$ с численным $\int c\, dz/H$, нуль $q$ на $\mu_{\text{cross}} = \sqrt{3-\sqrt{5}}$, формула возраста); \texttt{[1]} якорный SN-фит Pantheon+SH0ES (автоскачивание каталога и ковариаций, STAT+SYS, явный $M$) с $1\sigma$-погрешностями параметров из гессиана $\chi^2$; \texttt{[2]} контроль STAT-ONLY; \texttt{[3]} совместный SN+BAO с полной $12\times12$ ковариацией и восстановление линейки $r_d$; \texttt{[4]} $Om(z)$-диагностика со сканированием $H_0$ в двух режимах ковариации хронометров (диагональная и полная Moresco et al.\ 2020); \texttt{[5]} предсказания высокого $z$; \texttt{[6]} космическое время $t(z)$ DSIC vs $\Lambda$CDM.

Часть II: \texttt{[7]} самопроверка опорных данных Planck 2018 ($d_{\text{LSS}} = r^*/\theta^*$); \texttt{[8]} недобор ядра и фиксация $\mu_d$ по углу $\theta^*$ (заимствованный $r^*$ помечен явно); \texttt{[9]} контроль тождественного нуля поправки ниже порога и таблица вырождения «амплитуда $\leftrightarrow$ порог»; \texttt{[10]} след $\mu_d$ по $z$, исправленный горизонт частиц и флаг согласованности формулировки P6 (галактики JWST выше порога).

Часть III: \texttt{[11]} (рост структур): интегрирование стандартного уравнения линейного роста $\delta'' + (2 + d\ln H/d\ln a)\delta' - (3/2)\Omega_m(a)\delta = 0$ на фоне DSIC (GR-предел), вычисление $f\sigma_8(z)$ и сверка с компиляцией «Gold-2017» (18 точек, ковариация WiggleZ учтена); доклад результата через $S_8$ ввиду вырождения $\Omega_{m0}$--$\sigma_8$, плюс сравнение с $\Lambda$CDM на тех же точках.

\appendix

\section{Сводка обозначений}

\begin{center}
\footnotesize
\setlength{\tabcolsep}{4pt}
\begin{tabularx}{\textwidth}{@{}c l >{\raggedright\arraybackslash}X >{\raggedright\arraybackslash}X@{}}
\toprule
\textbf{слой} & \textbf{символ} & \textbf{значение} & \textbf{знак / диапазон} \\
\midrule
0 & $\varphi = \omega t$ & фаза потока (время) & монотонно растёт \\
0 & $\omega,\ \dot{\varphi}$ & скорость потока & $> 0$ \\
1 & $R$ & бюджет масштаба (норма) & $> 0$ \\
1 & $S = R\sin(\varphi/2)$ & масштаб пространства & со знаком (нечётн.) \\
1 & $O = R\cos(\varphi/2)$ & масштаб объектности & со знаком (чётн.) \\
1 & $\sigma_S, \sigma_O$ & знаковая фаза полюсов & $\pm 1$ \\
1 & $O_{\text{full}} = O\sigma_S$ & объектность с наследованным знаком & со знаком \\
1 & $W = -R^2/2$ & Вронскиан (квадратура) & const \\
2 & $P = |S| = R\sqrt{1-\mu^2}$ & наблюдаемый масштаб пространства & $\ge 0$ \\
2 & $Q = |O| = R\mu$ & наблюдаемый масштаб объектности & $\ge 0$ \\
2 & $\mu = Q/R = |\cos(\varphi/2)|$ & наблюдаемая фаза (доля объектности) & $\in [0,1]$, период $2\pi$ \\
3 & $L(\mu) \propto \mu$ & линейка наблюдателя & --- \\
3 & $\tau,\ d\tau = dt/\mu$ & собственное время наблюдателя & --- \\
3 & $c$ & скорость света & const \\
3 & $b = \sqrt{1-\mu^2}/\mu^2$ & наблюдаемый масштабный фактор & --- \\
3 & $z$ & красное смещение, $1+z = b(\mu_0)/b(\mu_e)$ & $\ge 0$ \\
3 & $s = \pm 1$ & направление ветви ($-1$ расшир., $+1$ сжатие) & --- \\
3 & $k = (2/\pi)\omega$ & скорость $\mu$ вдоль ветви & $> 0$ \\
3 & $H_\tau,\ q$ & темп расширения, параметр замедления & --- \\
3 & $w_{\text{eff}} = (2q-1)/3$ & эффективное уравнение состояния (перевод на язык ОТО, \S5.6) & --- \\
3 & $\chi, d_L, \text{DM}$ & сопутств., светимостное расст., модуль расст. & --- \\
3 & $\mu_0, \mu_e$ & фаза сегодня / при излучении & $\in [0,1]$; ветвь приёма $\mu_e \in (\mu_0, 1)$ \\
3 & $z_t$ & красное смещение перехода замедл.$\to$ускор. & $\approx 0.77$ \\
\bottomrule
\end{tabularx}
\end{center}

\medskip
\noindent\textbf{Обозначения Части II (второй этаж):}

\begin{center}
\footnotesize
\setlength{\tabcolsep}{4pt}
\begin{tabularx}{\textwidth}{@{}l >{\raggedright\arraybackslash}X >{\raggedright\arraybackslash}X@{}}
\toprule
\textbf{символ} & \textbf{значение} & \textbf{статус} \\
\midrule
связанность & вторая сущность второго этажа (в ядре отсутствует) & постулируется (P6) \\
$\mu_d = 0.980640$ & порог отлипания ($z_{\text{detach}} \approx 4.53$) & измеренная константа (калибровка по $d_{\text{LSS}}$) \\
$A \equiv 1$ & амплитуда поправки $1/P$ & фиксирована постулатом P6 \\
$D_M^{\text{испр}}$ & сопутствующее расстояние с добором $\arcsin$-поправки & --- \\
$d_{\text{LSS}}, z^*, r^*, \theta^*$ & опорные величины Planck 2018 & внешние данные \\
$r_{ij}, M_i, G$ & соотношение масштабов пары, массы, ньютоновская константа & модуль Прил.\ B, не из ядра \\
$g_{ij}, w_{ij}$ & связанность пары, вес распределения пространства & модуль Прил.\ B \\
$\kappa$ & темп перераспределения пространства & свободный параметр Прил.\ B \\
\bottomrule
\end{tabularx}
\end{center}

\section{Программа развития: гравитация и распределение пространства}

\begin{quote}
\textbf{Статус раздела.} Это не результат, а эскиз направления. Модуль ниже феноменологический: он использует ньютоновские величины $G, M_i$, которых в перетекании $S \leftrightarrow O$ нет; параметр $\kappa$ свободен; сверка с данными не проводилась. Он вынесен в приложение сознательно --- чтобы отделить проверенное ядро (Часть I) и откалиброванный второй этаж (Часть II) от гипотез, ещё не доведённых до количественной проверки. Ценность модуля --- архитектурная: он показывает, что связанность (Часть II) может иметь и позднее, гравитационное лицо, замыкая раннее (порог отлипания) и позднее (распределение пространства) на одну сущность.
\end{quote}

Позднее проявление той же связанности. Когда объекты падают в себя (\S3.4), гравитационно связанные группы падают согласованно --- у них проступает \textbf{общий центр падения}: чем сильнее объекты связаны, тем отчётливее выражен общий центр и тем меньше пространства высвобождается между ними; между слабо связанными --- больше. Петля самоусиливается: разлёт ослабляет связь $\to$ область получает больше пространства $\to$ разлёт ускоряется. Плотные структуры остаются компактными, разреженные получают больше проступающего пространства --- качественно это согласуется с наблюдаемым: пространство прибывает не там, где плотно, а там, где пусто. Пространство перераспределяется, глобальный баланс сохраняется, анизотропия локальна.

Обозначения модуля: $r_{ij}$ --- соотношение масштабов пары; $M_i$ --- масса объекта; $G = 1$; $\kappa$ --- свободный параметр (темп перераспределения).
\begin{align}
\text{связанность:} &\quad g_{ij} = \frac{G M_i M_j}{r_{ij}^2} \\
\text{градиент потенциала:} &\quad |\nabla\Phi_{ij}| = \left| \sum_k \frac{G M_k (\text{mid}_{ij} - \text{pos}_k)}{|\text{mid}_{ij} - \text{pos}_k|^3} \right| \\
\text{вес пространства:} &\quad w_{ij} = \frac{1/|\nabla\Phi_{ij}|}{Z}\,, \qquad Z = \sum_{i\neq j} \frac{1}{|\nabla\Phi_{ij}|} \\
\text{прирост:} &\quad \Delta r_{ij} = \frac{\kappa\,|\Delta P|\,w_{ij}}{M_i + M_j} \\
\text{эволюция:} &\quad r_{ij}(\varphi+\Delta\varphi) = r_{ij}(\varphi) + \Delta r_{ij} \quad (\text{только растёт}) \\
\text{сохранение:} &\quad \sum_{i\neq j} \Delta r_{ij}\,\frac{M_i + M_j}{\kappa} = |\Delta P(\varphi)|
\end{align}
Свойства весов: $w_{ij} \ge 0$, $\sum w_{ij} = 1$; пара в пустоте (плоское поле) получает большой вес, пара в кластере (крутое поле) --- малый.

Статус модуля --- феноменологический: он использует ньютоновские величины $G, M_i$, которых в перетекании $S \leftrightarrow O$ нет; $\kappa$ свободен; сверка с данными не проводилась. Но архитектурно он теперь стоит на том же основании, что и \S11: связанность --- одна сущность, у которой два лица. Раннее --- порог отлипания (условие, при котором связанность монолита разрывается и линейки наблюдателя обретают смысл); позднее --- распределение проступающего пространства в местах наименьшей гравитационной связанности. Программа-минимум: вывести $\mu_d$ как условие на связанность из того же корня, что и веса $w_{ij}$, --- тогда обе конструкции станут следствиями одной величины, а не двумя довесками.

\bigskip
\noindent\textit{Спасибо, что дочитали мою захватывающую статью.}

\medskip
\noindent Авантюрист --- А. Засеев

\end{document}
